\chapter{Presentación de la empresa}

Synaptix presenta su actividad, organización y forma de gobierno para el proyecto Puerto Deportivo Bahía Panitao. La experiencia y los antecedentes habilitantes se documentan en los formularios de la oferta; este capítulo describe las capacidades y responsabilidades que sustentan los subdocumentos técnicos siguientes.

\section{Presentación de la empresa}

Synaptix Ingeniería de Software desarrolla plataformas operacionales para organizaciones que administran activos físicos. Su oferta comprende arquitectura de software, integración entre sistemas y equipos de terreno, gestión de datos, trazabilidad y servicios de operación tecnológica. Para Bahía Panitao, estas disciplinas se aplican a una plataforma híbrida que debe mantener las funciones críticas del recinto durante interrupciones de conectividad.

La identificación legal, financiera y de representación de la empresa se incorpora en el Sobre N\textsuperscript{o} 1 de la oferta, conforme a las Bases Administrativas \parencite[arts.~34.1 y 39, pp.~22 y 25]{bases-admin}. \textbf{[PENDIENTE: razón social inscrita, RUT, domicilio, representante legal y documentos de respaldo.]}

\section{Estructura organizacional}

La Figura~\ref{fig:organizacion-synaptix} muestra las funciones que Synaptix articula para la prestación propuesta. Dirección coordina el contrato; Arquitectura e Ingeniería diseña e integra la solución; Operaciones sostiene el servicio; Datos y Seguridad establece controles transversales; Calidad verifica entregables y evidencia. La dotación nominal del proyecto se presenta por separado en el Formulario T-8.

\begin{figure}[htbp]
  \centering
  \begin{tikzpicture}[
    area/.style={draw=SynNavyLight,rounded corners=2mm,fill=SynSurface,text=SynInk,align=center,font=\small,minimum height=12mm,text width=42mm},
    direccion/.style={area,fill=SynNavy,text=SynWhite,text width=50mm},
    conexion/.style={draw=SynNavyLight,thick,-{Latex}}
  ]
    \node[direccion] (dir) {Dirección general\\y gestión contractual};
    \node[area] (arq) at ([xshift=-42mm,yshift=-22mm]dir.south) {Arquitectura e Ingeniería\\diseño e integración};
    \node[area] (ops) at ([xshift=42mm,yshift=-22mm]dir.south) {Operaciones y Servicios Gestionados\\continuidad y soporte};
    \node[area] (datos) at ([xshift=-42mm,yshift=-50mm]dir.south) {Datos y Seguridad\\gobierno y controles};
    \node[area] (calidad) at ([xshift=42mm,yshift=-50mm]dir.south) {Calidad y Conocimiento\\verificación y transferencia};
    \draw[conexion] (dir.south west) -- (arq.north);
    \draw[conexion] (dir.south east) -- (ops.north);
    \draw[conexion] (arq.south) -- (datos.north);
    \draw[conexion] (ops.south) -- (calidad.north);
    \draw[SynCyan,thick,dashed] (datos.east) -- node[above,font=\scriptsize,text=SynMuted]{control transversal} (calidad.west);
  \end{tikzpicture}
  \caption{Funciones de Synaptix para la prestación propuesta}
  \label{fig:organizacion-synaptix}
  \FuenteFigura{Elaboración de Synaptix.}
\end{figure}

Arquitectura e Ingeniería reúne el diseño de aplicaciones, interfaces y componentes de borde. Operaciones administra despliegues, observabilidad, respaldo, atención y continuidad. Datos y Seguridad define reglas para acceso, conservación, intercambio y protección de información. Calidad y Conocimiento revisa la trazabilidad de requisitos, registra no conformidades y prepara la transferencia al cliente. Dirección resuelve las decisiones que afectan alcance, plazo, costo o compromisos contractuales. \textbf{[PENDIENTE: dotación por unidad y responsable titular de cada función.]}

\section{Gobierno interno de calidad, seguridad y conocimiento}

Para este proyecto, Synaptix establecerá una línea base de requisitos y decisiones técnicas. La jefatura del proyecto aprobará cambios de alcance; Arquitectura mantendrá el registro de decisiones y revisará las interfaces; Calidad contrastará entregables con requisitos y criterios de aceptación antes de cada hito. Las actas de revisión, la matriz de trazabilidad y el registro de incidencias serán las evidencias de aplicación de este control.

Seguridad definirá perfiles de acceso, revisión de privilegios, tratamiento de incidentes y controles para el trabajo sin conexión. Las excepciones se registrarán con responsable, plazo y medida compensatoria. Operaciones revisará eventos, respaldos y pruebas de recuperación; las conclusiones se incorporarán al registro de riesgos y al plan de mejora del servicio.

El conocimiento técnico se conservará en documentación versionada: arquitectura, procedimientos de operación, inventario de integraciones y manuales de uso. Antes del traspaso de cada etapa, Calidad comprobará que esos documentos correspondan a la versión desplegada y que la contraparte pueda ejecutar los procedimientos acordados. \textbf{[PENDIENTE: responsables nominales, periodicidad aprobada de las revisiones y repositorios de evidencia.]}

\section{Experiencia y certificaciones}

Synaptix presentará su experiencia comparable mediante el Formulario T-6 independiente, con tres proyectos finalizados y en operación dentro del plazo exigido. La ficha debe incluir cliente, fechas, alcance, arquitectura, nivel de servicio, volumen, rol y contraparte de referencia. Al menos un proyecto debe demostrar arquitectura híbrida y otro, disponibilidad comprometida igual o superior a 99,5\,\% \parencite[art.~34.1, p.~22; Formulario T-6, p.~56]{bases-admin}. \textbf{[PENDIENTE: completar las tres fichas T-6 con antecedentes comprobables y cartas de referencia.]}

La certificación institucional ISO/IEC 27001 se presenta mediante certificado vigente o plan formal de certificación con hitos verificables, según la alternativa prevista en las Bases \parencite[art.~34.1, p.~22]{bases-admin}. Cualquier otra certificación institucional o individual se informará con su titular, alcance, vigencia y medio de verificación. \textbf{[PENDIENTE: certificados institucionales, plan formal cuando proceda y nómina de certificaciones individuales.]}

Los estados financieros, la capacidad de constituir garantías y los certificados laborales corresponden a la documentación administrativa de la oferta \parencite[art.~34.1, p.~22]{bases-admin}. \textbf{[PENDIENTE: incorporar los antecedentes exigidos en el Sobre N\textsuperscript{o} 1.]}

\section{Estructura para el proyecto}

La organización del proyecto separa dirección contractual, construcción de la plataforma y operación del servicio. La jefatura de proyecto concentra la relación formal con Bahía Panitao y coordina hitos, presupuesto y cambios. Arquitectura de Solución define los componentes y sus interfaces; los equipos de desarrollo, datos e integración implementan y prueban las funciones. Seguridad revisa los controles de acceso y tratamiento de datos. Calidad verifica los criterios de aceptación y Operaciones prepara el paso a producción y el servicio continuo.

El Formulario T-8 individualiza a las personas que cubrirán los roles clave, su dedicación y sus compromisos. Este organigrama funcional no sustituye esa nómina. \textbf{[PENDIENTE: asignación nominal, dedicaciones, currículos y cartas de compromiso del equipo clave en T-8.]}

La coordinación con la marina se realizará mediante una contraparte técnica para decisiones de alcance y aceptación, y responsables operacionales para validar procesos de muelle, escuela, varadero y administración. Los mecanismos de implantación y la dotación por etapa se desarrollan en los subdocumentos de solución y planificación.

\section{Alianzas empresariales}

La propuesta contempla colaboración con el proveedor de nube, fabricantes o integradores de instrumentación y prestadores de servicios especializados cuando su participación sea necesaria para cumplir el alcance. La condición de socio del proveedor de nube, o el acuerdo formal con un socio certificado participante, se acredita conforme al artículo 34.1 de las Bases Administrativas \parencite[p.~22]{bases-admin}. \textbf{[PENDIENTE: identificar entidades participantes, función contractual y certificados o cartas de compromiso.]}

Cada alianza incorporada a la oferta tendrá un responsable de integración, entregables definidos y un mecanismo de coordinación con Synaptix. La descripción de productos y componentes específicos corresponde al subdocumento de arquitectura y al Formulario T-11; aquí se establece la responsabilidad empresarial por su integración.
